\documentclass[conference]{IEEEtran}
\usepackage{cite}
\usepackage{amsmath,amssymb,amsfonts}
\usepackage{graphicx}
\usepackage{textcomp}
\usepackage{xcolor}
\usepackage{booktabs}
\usepackage{hyperref}

\begin{document}

\title{Proactive Kubernetes Autoscaling Using Patch-Based Transformer Models for Latency-Sensitive Microservices}

\author{\IEEEauthorblockN{Dhruv}
\IEEEauthorblockA{\textit{Department of Computer Science} \\
\textit{Indian Institute of Technology Delhi}\\
New Delhi, India \\
dhruv@iitd.ac.in}
}

\maketitle

\begin{abstract}
The Kubernetes Horizontal Pod Autoscaler (HPA) follows a reactive control loop that scales resources only after threshold violations are observed, introducing a systemic lag of 60--90 seconds during sudden traffic surges. This delay is particularly damaging in latency-sensitive microservice deployments, where even brief periods of under-provisioning degrade user experience and violate service-level objectives. We present a proactive autoscaling framework that replaces the reactive feedback loop with a predictive controller built on PatchTST, a patch-based Transformer architecture for time-series forecasting. Our evaluation spans two heterogeneous microservice benchmarks---the Google Online Boutique (10 services, Go/Python) and Weaveworks Sock Shop (13+ services, Java/Node.js/Go)---deployed on Google Kubernetes Engine clusters with Prometheus-based observability at 2--15 second resolution. Over 117 hours of continuous traffic simulation (28,108 temporally aligned observations for Online Boutique and 7,147 for Sock Shop), we demonstrate that PatchTST achieves an RMSE of 0.0179 on Online Boutique and 0.0917 on Sock Shop, representing a 97.7\% and 85.5\% reduction in mean absolute error compared to the reactive HPA baseline, respectively. The proactive controller detects 98.2\% of traffic spikes on Online Boutique and 89.5\% on Sock Shop, while maintaining false positive rates below 5\%---compared to HPA's 52\% recall and 43--48\% resource waste under identical conditions. A zero-shot transfer learning experiment, in which a model trained exclusively on Online Boutique data is evaluated on Sock Shop without fine-tuning, yields a 69\% improvement over HPA, confirming that the learned representations generalize across architecturally distinct microservice platforms. These results establish patch-based Transformer forecasting as a robust and generalizable alternative to threshold-driven autoscaling in production Kubernetes environments.
\end{abstract}

\begin{IEEEkeywords}
Kubernetes, Autoscaling, Transformer, PatchTST, Time-Series Forecasting, Cloud Engineering, Microservices, Transfer Learning
\end{IEEEkeywords}

\end{document}
